\documentclass{article}
\begin{document}
\section{Many floats}
Figures \ref{f1}, \ref{f2}, \ref{f3} and Table~\ref{t1} come later.

\begin{figure}[t]\centering\mbox{One}\vspace{150pt}\caption{First.}\label{f1}\end{figure}
\begin{figure}[t]\centering\mbox{Two}\vspace{150pt}\caption{Second.}\label{f2}\end{figure}
\begin{figure}[htbp]\centering\mbox{Three}\vspace{150pt}\caption{Third.}\label{f3}\end{figure}
\begin{table}[p]\centering\begin{tabular}{|c|c|}\hline x & y\\\hline\end{tabular}\caption{Page only.}\label{t1}\end{table}

A paragraph of text after the floats that is long enough to go over a line or two so
that the page has some text on it before the next floats arrive in the source text.

\begin{figure}[!h]\centering\mbox{Here!}\caption{Here if possible.}\end{figure}

More text, and a footnote\footnote{A note on the same page as floats.} to mix inserts
with floats on one page.

\begin{figure}[b]\centering\mbox{Bottom}\caption{Bottom one.}\end{figure}
\begin{table}[tb]\centering\begin{tabular}{l}a\\b\end{tabular}\caption{Top or bottom.}\end{table}

Closing text.
\begin{figure}\centering\mbox{Default}\vspace{300pt}\caption{Default placement, tall.}\end{figure}
\begin{figure}[h]\mbox{Late}\caption{Left at the end.}\end{figure}
\end{document}
